% 04b-identity.tex — v3 §3.2.
\subsection{The margin identity and the push-through form}\label{sec:noise-identity}
% v3: §3.2

\paragraph{Bookkeeping} (Lemmas~\ref{lem:count}, \ref{lem:excess}, Corollary~\ref{cor:lag}). Rigid integers and
elements of $\Abs$ sit on their own unit jumps of $T$ (jump slots). Every level crossing of $T$ is either a jump slot or a
continuous slot, so on every window $W$
\[
  \#\{\text{continuous slots in }W\}=\int_WT'_{\mathrm{cont}}\pm2 .
\]
% Logged correction 95 (CORRECTIONS-vs-v3.md): v3 states the bookkeeping on slot windows with a +-2 for "the two slot
% boundaries", but the margin identity applies it to the physical window W_y, whose end cells are cut, and its nu has no edge
% term (Lemma noise adds one). On slot windows the identity is exact; for any interval the residual E_edge of the two
% partial end cells is defined and bounded below, and it now enters nu.
By Lemma~\ref{lem:count}, on a slot window $W=(k,j]$ (the slots $k+1,\dots,j$, with physical interval $[s_{k+1},s_{j+1})$) the
number of primes decided equals the number of continuous slots, minus the non-absorbed jittered composites landing in $W$,
plus the change in the queue:
\begin{equation}\label{eq:bookkeeping}
  \sigma(W)=\#\{\text{continuous slots in }W\}-C_{\mathrm{jit}}(W)+\Delta Q(W),\qquad\Delta Q(W):=Q_j-Q_k .
\end{equation}
Indeed $\sum_{i=k+1}^j(P_i+A_i-1)=Q_j-Q_k$ by the greedy rule, and the composites with value in $[s_{k+1},s_{j+1})$, counted by
$\sum_{i=k+1}^jA_i$, are the elements of the jump slots of $W$, one each, and the $C_{\mathrm{jit}}(W)$ others. Here $C_{\mathrm{jit}}$ counts the composites $n\notin\Abs$ whose
latest prime $q$ (in decision order) is jittered; the cofactor may be any element of $\NN$. Write
$C_{\mathrm{jit}}=C^{(1)}+C^{(\ge2)}$, according to whether $q\,\|\,n$ or $q^2\mid n$.

\emph{Physical windows.} For an arbitrary interval $W$, let $k$ be the last slot that starts before $W$ ($k:=j_0-1$ if there
is none), and let $k+1,\dots,j$ ($j\ge k$) be the slots that start in $W$. Then $\sigma(W)$ and the continuous slots in $W$
refer to these slots (\S\ref{sec:greedy}), while $C_{\mathrm{jit}}(W)$ counts by value. Put $\Delta Q(W):=Q_j-Q_k$, and let
$W^\circ:=[s_{k+1},s_{j+1})$ be the physical interval of the slot window $(k,j]$. The \emph{edge residual} of $W$ is
\begin{equation}\label{eq:edge}
\begin{aligned}
  E_{\mathrm{edge}}(W)&:=\sigma(W)-\#\{\text{continuous slots in }W\}+C_{\mathrm{jit}}(W)-\Delta Q(W)\\
  &=C_{\mathrm{jit}}(W\setminus W^\circ)-C_{\mathrm{jit}}(W^\circ\setminus W),
\end{aligned}
\end{equation}
where the second form is \eqref{eq:bookkeeping} for $(k,j]$. Here $W\setminus W^\circ$ lies in the cell $J_k$ (for $k=j_0-1$,
below $s_{j_0}$, where there are no jittered composites, Lemma~\ref{lem:count}), and $W^\circ\setminus W$ (empty if $j=k$) lies
in the cell $J_j$, which ends after $W$. These are the two partial end cells. Since $\sigma(W)$, the continuous-slot count and
$\Delta Q(W)$ refer to the same slots, the slot starts cancel, and only the composites in the cut parts of the two end cells
remain. As $A_i\le Q_i+1$ (by $Q_i=Q_{i-1}+A_i+P_i-1$ with $Q_{i-1},P_i\ge0$),
\[
  -A_j\le E_{\mathrm{edge}}(W)\le A_k,\qquad|E_{\mathrm{edge}}(W)|\le1+\max(Q_k,Q_j),
\]
with $A_{j_0-1}:=0$. If $W$ is the physical interval of a slot window, then $W^\circ=W$ and $E_{\mathrm{edge}}(W)=0$. (Its $k,j$ may differ from those of the slot window by an empty slot at an end; $W^\circ$ is the same.)
\begin{itemize}
\item $\Comp$ is the predictable compensator of $C^{(1)}$. Given $\mathcal F_{q-}$, the element $q\cdot w$ ($q\nmid w$)
  lands in $W$ with conditional probability $|J_q\cap W/w|/H(s_q)$, and $q\cdot w$ is formed only if $q$ is decided
  after every prime of $w$, i.e.\ $s_q>\tau_w$. Hence
  \[
    \Comp(W)=\sum_{w\ge2}\int_{W/w}J(\Pi_{\tau_w}\sigma)\cdot\mathbf 1[qw\notin\Abs],
  \]
  where $\Pi_\tau\sigma$ keeps the slots $s_q>\tau$.
% Logged correction 83 (CORRECTIONS-vs-v3.md): the sum defining Comp^(2) is over the jittered q (as is C^(>=2)), the m decided
% before q, and the elements outside Abs; v3 leaves the range of q, the latest-prime filter and the absorption indicator implicit.
\item $\Comp^{(2)}$ is the predictable compensator of $C^{(\ge2)}$:
  $\Comp^{(2)}(W):=\sum_q\sum_{a\ge2}\sum_m\Prob(q^am\in W,\ q^am\notin\Abs\mid\mathcal F_{q-})$, over the jittered primes $q$
  and the $m\in\NN$ with $q\nmid m$ and $\tau_m<s_q$ (including $m=1$).
\item $\Mart':=(C^{(1)}-\Comp)+(C^{(\ge2)}-\Comp^{(2)})=:\Mart+\Mart^{(2)}$. Both parts are martingales in decision
  order. They are controlled by separate Freedman bounds (\S\ref{sec:noise-perscale}), because the increment of $q$ enters
  $\Mart$ through $W/q$ and $\Mart^{(2)}$ through $W/q^a$.
\end{itemize}

Since $\int_{W_x/w}g=\frac1w\int_{W_x}g(y/w)\,dy$,
\[
  \Comp(W_x)/H(x)=(B\Kpt\rhot)(x)-(BK_{\Abs}\rhot)(x)+\NT(x).
\]
\begin{itemize}
\item $K_{\Abs}$ is the part of $\Kpt$ whose landing is absorbed ($y/w<y(y)$, $\log y\ge L_*$).
\item The cutoff $\mathbf 1[y/w\ge\Xm]$ in $\Kpt$ costs nothing, since $\rhot$ vanishes below $\Xm$.
\item $\NT$ replaces the slot filter $s_q>\tau_w$ by the landing filter $y/w>\tau_w$ on the non-absorbed part. It is
  supported on \emph{$\tau$-near-threshold} cofactors, i.e.\ those with $|\tau_w-y/w|\le2H(y/w)$, and
  \begin{equation}\label{eq:NT}
    |\NT(y)|\le2\eta(y)\cdot\sup\rhot,\qquad\eta(y):=\sum_{w\ \tau\text{-near at }y}\frac1w .
  \end{equation}
\item $\NT_{\mathrm{pt}}$ denotes $\NT$ before the $B$-average, and $m:=T'-\Kpt\rhot-\NT_{\mathrm{pt}}$ is the
  \emph{pointwise margin density}.
\end{itemize}

\begin{proposition}[the margin identity]\label{prop:margin-identity}
% v3: Identity (3.1)
% Logged correction 95: nu contains the edge residual of W_y (eq:edge) and the box term E_box, so the identity is exact.
For every $y\ge X_0$,
\begin{equation}\label{eq:margin-identity}
\begin{gathered}
  \rhot+B\Kpt\rhot=F,\qquad F:=\tilde T'-\nu-\NT,\\
  \nu(y):=\frac{\Mart'(W_y)-\Delta Q(W_y)-E_{\mathrm{edge}}(W_y)+\Comp^{(2)}(W_y)+E_{\mathrm{box}}(y)}{H(y)},
\end{gathered}
\end{equation}
where $\Delta Q(W_y)$ and $E_{\mathrm{edge}}(W_y)$ are those of the physical window $W_y$ \eqref{eq:edge}, and
$E_{\mathrm{box}}(y):=\sigma(W_y)-H(y)\rhot(y)$. Moreover $|E_{\mathrm{box}}|\le2.01$.
\end{proposition}

\begin{proof}
% Logged correction 95: v3 applies the slot-window bookkeeping to W_y and puts the box term into theta, although tilde T' is
% defined without it; both residuals are now explicit terms of nu.
By \eqref{eq:edge} for the physical window $W_y$, exactly,
\[
  \sigma(W_y)=\#\{\text{continuous slots in }W_y\}-C_{\mathrm{jit}}(W_y)+\Delta Q(W_y)+E_{\mathrm{edge}}(W_y).
\]
Divide by $H(y)$. On the left,
$\sigma(W_y)/H(y)=\rhot(y)+E_{\mathrm{box}}(y)/H(y)$. On the right, write $C_{\mathrm{jit}}=\Comp+\Comp^{(2)}+\Mart'$, use the display
for $\Comp(W_y)/H(y)$ above, and move $B\Kpt\rhot$ to the left. The terms $BK_{\Abs}\rhot$ and
$\#\{\text{continuous slots in }W_y\}/H$ make up $\tilde T'$, the term $-\NT$ remains, and the rest is $-\nu$; there is no error term.

The box term: $H(y)\rhot(y)=\sum_qH(y)H(s_q)^{-1}\mathbf 1[y-H(s_q)/2<s_q\le y+H(s_q)/2]$, while $\sigma(W_y)$ counts the $q$ with
$s_q\in W_y$. Every box is centred, with width $H(s_q)=H(y)(1+O(\eps H(y)/y))$ for these $q$ by \eqref{eq:RH}, and distinct
decided primes have slot starts at least $1-10^{-8}$ apart (they are continuous-slot starts, \S\ref{sec:greedy}). So at most one prime near each end of
$W_y$ is counted differently, and the weights $H(y)/H(s_q)$ change the rest by $O(\eps H(y)^2/y)\le0.01$. Hence
$|E_{\mathrm{box}}|\le2.01$.
\end{proof}

The compensator margin density at $y$ is
\begin{equation}\label{eq:M}
  M(y):=\tilde T'(y)-(B\Kpt\rhot)(y)=\rhot(y)+\nu(y)+\NT(y).
\end{equation}
The noise $\nu$ enters only at the current scale $y$: it depends on the data on $W_y$ and on the two slot cells at its ends (Lemma~\ref{lem:noise}). Every other
quantity in \eqref{eq:margin-identity} at $y$ involves $\sigma$ only at scales $\le y/2+H$.

\paragraph{The push-through form.} $\Kpt$ maps functions supported on scales $\ge v$ to functions supported on scales
$\ge2v$. So $1+B\Kpt$ and $1+\Kpt B$ are invertible, and their Neumann series terminate on bounded ranges. For any such
operators $B,K$,
\[
  (1+BK)^{-1}=(1-B)+B(1+KB)^{-1}\big(1-K(1-B)\big),
\]
since $(1+BK)(1-B)+B(1-K(1-B))=1-B+BK-BKB+B-BK+BKB=1$. With $K=\Kpt$ and $\Rres:=(1+\Kpt B)^{-1}$, \eqref{eq:margin-identity}
gives
\begin{equation}\label{eq:pushthrough}
  \rhot=(1-B)F+B\Rres F-B\Rres\Kpt(1-B)F .
\end{equation}
\begin{itemize}
\item $\Rres$ is the \emph{long} operator: at each step the sharp landing filter, followed by an $H(u)$-average at the
  landing scale $u$. Its kernels $C_n$ are studied in \S\ref{sec:kernels}.
\item $(1-B)$ acts only on the source. The response $\rhot$ is never differentiated, so no regularity lemma for it is needed.
\item Applying the pointwise inverse $(1+\Kpt)^{-1}$ to \eqref{eq:margin-identity} instead would leave the response term
  $(1-B)\Kpt\rhot$, which is of size $\asymp1/\log y$ in sup norm. This is why the push-through form is used.
\end{itemize}
Note also the push-through identity $(1+B\Kpt)^{-1}B=B\Rres$ and its consequence $\Kpt B\Rres=1-\Rres$.

\paragraph{The jump split.} Split $T'=\Tc+J$ with $J=\jz\mathbf 1_{[X_0,\infty)}$ (\S\ref{sec:noise-notation}). The
$J$-part of $F$ is $BJ$. By the push-through identity its contribution to $\rhot$, and hence to the margin
$M=\rhot+\nu+\NT$, is exactly $B\Rres J$. For the rest put
\begin{equation}\label{eq:decomp}
  F_c:=F-BJ=\Tc+F_r,\qquad F_r:=(B-1)\Tc+\vartheta/H-\nu-\NT,\qquad G:=\Kpt(1-B)F_c .
\end{equation}
By \eqref{eq:pushthrough} applied to $F_c$,
$M=B\Rres J+(1-B)F_c+B\Rres F_c-B\Rres G+\nu+\NT$.

\emph{Placement convention.} The absorbed source $\sigma_{\Abs}$ is part of $\Tc$, never of $F_r$. Only its
$H$-oscillation enters, through $(B-1)\Tc$ and $o_u$ (see \eqref{eq:osc}). The jump enters only through $B\Rres J$, that is,
through Lemma~\ref{lem:xb}.
